% --- ARCHIVO: Anexos_Conceptos/anexo_cap7.tex ---
% VERSION 2 — Entradas añadidas:
%   · \label{anexo:ufls}  — referenciado en §7.1.1, no existía
%   · \label{anexo:pu}    — referenciado en §7.1.2, no existía
% Correcciones:
%   · El capítulo usa \ref{anexo:statcom} → corregido a \ref{anexo:syncon}
%     en resiliencia_futuro_v4.tex; esta entrada permanece bajo el label
%     correcto \label{anexo:syncon}
%   · El capítulo usaba \ref{anexo:grid_forming} → corregido a
%     \ref{anexo:grid_following} en resiliencia_futuro_v4.tex
% Tabla de referencias cruzadas al final del archivo (comentario)

\chapter{Conceptos técnicos y regulatorios del Capítulo~7}
\label{anexo:conceptos_cap7}

% ------------------------------------------------------------------
\section{Coste Nivelado de la Energía (LCOE)}
\label{anexo:lcoe}
% ------------------------------------------------------------------
El \textit{Coste Nivelado de la Energía} (LCOE) es la métrica económica
estándar que compara el coste unitario de producción entre distintas
tecnologías a lo largo de su vida útil. Su principal limitación sistémica
es que ignora el valor de los servicios ancilares aportados a la red;
así, aunque los IBR presentan un LCOE tendente a cero, no proveen de
forma nativa la inercia o potencia de cortocircuito que las centrales
síncronas tradicionales ---con mayor LCOE--- aportaban implícitamente,
originando la paradoja de escasez de firmeza en sistemas altamente
renovables.

% ------------------------------------------------------------------
\section{Servicios Esenciales de Confiabilidad (ERS)}
\label{anexo:ers}
% ------------------------------------------------------------------
Los \textit{Servicios Esenciales de Confiabilidad} (ERS) agrupan los
atributos físicos indispensables para la operación segura de la red,
tales como la inercia, la potencia de cortocircuito, la respuesta rápida
de frecuencia y el control dinámico de tensión. Dado que las tecnologías
asíncronas no proveen estos servicios de forma inherente como lo hacían
las máquinas rotativas convencionales, los nuevos marcos regulatorios
(como la reforma del P.O.~7.4 o el NC~RfG~2.0) exigen la creación de
mercados específicos que definan y remuneren explícitamente su provisión.

% ------------------------------------------------------------------
\section{BESS con inversores \textit{Grid-Forming} (BESS-GFM)}
\label{anexo:bess_gfm}
% ------------------------------------------------------------------
Los \textit{Sistemas de Almacenamiento en Baterías con Inversores
Formadores de Red} (BESS-GFM) combinan alta densidad electroquímica con
electrónica de potencia capaz de operar como una fuente de tensión ideal
y autónoma. A diferencia de los inversores seguidores de red
(\textit{grid-following}), la topología GFM permite al sistema ofrecer
arranque autónomo (\textit{Black Start}), inercia sintética, respuesta
rápida de frecuencia (FFR) y soporte activo de tensión, asumiendo así
las responsabilidades de estabilización sistémica críticas en redes
descarbonizadas.

% ------------------------------------------------------------------
\section{Fast Frequency Response (FFR)}
\label{anexo:ffr}
% ------------------------------------------------------------------
La \textit{Respuesta Rápida de Frecuencia} (FFR) es un servicio de
estabilización subcíclica, diseñado para sistemas de electrónica de
potencia, que inyecta un bloque masivo de potencia activa en la ventana
temporal crítica (típicamente inferior a 0{,}25~s) previa a la actuación
de los reguladores mecánicos tradicionales. Al mitigar drásticamente la
tasa de caída de frecuencia (RoCoF), la FFR compensa eficazmente la falta
de inercia rotacional, habiendo demostrado en mercados pioneros como
ERCOT su capacidad para reducir los umbrales críticos de seguridad del
sistema.

% ------------------------------------------------------------------
\section{Estrategia \textit{Brownfield}}
\label{anexo:brownfield}
% ------------------------------------------------------------------
En ingeniería de infraestructuras energéticas, la estrategia
\textit{Brownfield} consiste en la reconversión de instalaciones
industriales existentes ---como las centrales térmicas o nucleares
clausuradas--- para dotarlas de nuevas funciones sistémicas. En este
contexto, implica conservar los grandes alternadores originales operando
en vacío como compensadores síncronos, aportando inercia natural y
potencia de cortocircuito, y aprovechando las subestaciones y líneas de
evacuación ya construidas para reducir drásticamente costes y tiempos de
implementación.

% ------------------------------------------------------------------
\section{Programa DS3 de EirGrid}
\label{anexo:ds3}
% ------------------------------------------------------------------
El programa \textit{Delivering a Secure, Sustainable Electricity System}
(DS3) es el marco pionero de servicios ancilares de EirGrid (Irlanda),
diseñado para operar el sistema insular con penetraciones renovables
asíncronas de hasta el 75~\%. El programa fragmenta y remunera la
estabilidad en 14~productos técnicos muy especializados ---como la
Respuesta Inercial Síncrona (SIR) indexada al mínimo técnico, o la FFR
con penalizaciones por rebote---, marcando el estándar internacional sobre
cómo traducir las necesidades físicas de una red descarbonizada a
incentivos económicos.

% ------------------------------------------------------------------
\section{Vehicle-to-Grid (V2G)}
\label{anexo:v2g}
% ------------------------------------------------------------------
La tecnología \textit{Vehicle-to-Grid} (V2G) habilita la bidireccionalidad
de las baterías de los vehículos eléctricos, permitiéndoles inyectar
potencia activa y reactiva hacia la red. Bajo el nuevo marco europeo
NC~RfG~2.0, las infraestructuras de recarga que agreguen más de 1~MW de
capacidad conjunta serán catalogadas como Módulos de Almacenamiento
significativos (Tipo~B), imponiéndoles la obligación de incorporar
capacidades \textit{Grid-Forming} para transformar estos megaparques de
movilidad en nodos activos de estabilización dinámica.

% ------------------------------------------------------------------
\section{Tasa de Cambio de Frecuencia (RoCoF)}
\label{anexo:rocof}
% ------------------------------------------------------------------
La \textit{Rate of Change of Frequency} (RoCoF, tasa de cambio de
frecuencia) cuantifica la velocidad de variación de la frecuencia del
sistema ante una perturbación, expresada típicamente en Hz/s. Es el
parámetro dinámico más crítico para la estabilidad transitoria: un RoCoF
elevado reduce el tiempo disponible para que los sistemas de regulación
actúen, acelerando la cascada de desconexiones de protecciones. El
colapso del 28A evidenció que en sistemas con baja inercia y débil
potencia de cortocircuito, un RoCoF de apenas 1--2~Hz/s es capaz de
superar los tiempos mecánicos de respuesta de los relés de deslastre de
carga (UFLS).

% ------------------------------------------------------------------
\section{Relés de Deslastre de Carga (UFLS)}
\label{anexo:ufls}
% ------------------------------------------------------------------
El \textit{Under-Frequency Load Shedding} (UFLS, deslastre automático de
carga por baja frecuencia) es el mecanismo de último recurso del sistema
de defensa: cuando la frecuencia cae por debajo de umbrales predefinidos
(habitualmente escalonados entre 49{,}0~Hz y 48{,}0~Hz en el Área
Síncrona de Europa Continental), los relés de UFLS desconectan cargas de
forma automática para restaurar el equilibrio generación-demanda. Su
limitación intrínseca es la velocidad de actuación: los interruptores
mecánicos de alta tensión requieren entre 60~ms y 100~ms para abrir, lo
que los hace inoperables como primera barrera ante perturbaciones con
RoCoF superior a 2~Hz/s. En sistemas con baja inercia ($H < 3$~s), el
umbral de 49~Hz puede alcanzarse en menos de 200~ms, comprimiendo la
ventana de actuación hasta hacerla mecánicamente inviable.

% ------------------------------------------------------------------
\section{Sistema por Unidad (p.u.)}
\label{anexo:pu}
% ------------------------------------------------------------------
El \textit{sistema por unidad} (p.u.) es una convención de normalización
utilizada en ingeniería eléctrica de potencia que expresa las magnitudes
del sistema (tensión, corriente, potencia, impedancia) como cocientes
adimensionales respecto a valores base de referencia. La base de tensión
suele tomarse como el valor nominal de la red en el nudo de análisis, y
la base de potencia como la potencia aparente nominal del equipo o del
sistema. La ventaja principal es la eliminación de las transformaciones
de escala al analizar redes con múltiples niveles de tensión
interconectados mediante transformadores. En el contexto del análisis de
inversores, la expresión de las corrientes de falta en p.u. permite
comparar directamente la capacidad de inyección de los inversores
(1{,}1--1{,}2~p.u.) con la de los generadores síncronos (5--7~p.u.)
con independencia de la potencia nominal de cada tecnología.

% ------------------------------------------------------------------
\section{\textit{Grid-Following} vs.\ \textit{Grid-Forming}: arquitecturas
de control de inversores}
\label{anexo:grid_following}
% ------------------------------------------------------------------
La topología \textit{Grid-Following} (GFL) modela al inversor como una
fuente de corriente controlada que depende de una medición externa de la
tensión de red (a través del lazo de seguimiento de fase, PLL). Su
ventaja es la simplicidad y el bajo coste; su limitación crítica es que
no puede operar de forma autónoma ni establecer tensión en redes débiles.
La topología \textit{Grid-Forming} (GFM) modela el inversor como una
fuente de tensión ideal detrás de una reactancia virtual, permitiendo
operación autónoma, inyección de corrientes de falta robustas e inercia
sintética. El NC~RfG~2.0 establece la transición hacia GFM como
obligatoria para nuevas instalaciones significativas.

% ------------------------------------------------------------------
\section{Inercia Sintética}
\label{anexo:inercia_sintetica}
% ------------------------------------------------------------------
La \textit{inercia sintética} es un algoritmo de control implementado en
inversores GFM que emula matemáticamente el comportamiento de la ecuación
de oscilación de un rotor electromecánico. El algoritmo mide
continuamente la derivada temporal de la frecuencia ($df/dt$) y ajusta
la potencia inyectada de forma proporcional, liberando o absorbiendo
energía del banco de baterías en decenas de milisegundos. Esta respuesta
subcíclica ralentiza la tasa de caída de frecuencia, ampliando el margen
temporal para la actuación de los sistemas de regulación convencionales,
y es equivalente funcionalmente a la masa cinética de las máquinas
síncronas tradicionales pero con mucha mayor velocidad de respuesta.

% ------------------------------------------------------------------
\section{Potencia de Cortocircuito}
\label{anexo:potencia_cortocircuito}
% ------------------------------------------------------------------
La \textit{Potencia de Cortocircuito} ($S_{sc}$) en un nudo es la
magnitud instantánea de corriente que el sistema puede inyectar ante una
falta de tensión. Define la rigidez eléctrica del nudo: un $S_{sc}$
elevado permite que las protecciones de distancia operen correctamente,
que las protecciones de sobrecorriente se coordinen selectivamente y que
los inversores mantengan sincronismo de sus algoritmos de control. Los
inversores GFM pueden contribuir a $S_{sc}$ mediante la inyección de
corrientes de falta controladas; sin embargo, su límite térmico
(1{,}1--1{,}2~p.u.) es mucho menor que el de las máquinas síncronas
(5--7~p.u.), requiriendo el complemento de compensadores síncronos en
redes débiles.

% ------------------------------------------------------------------
\section{Compensadores Síncronos (SynCons)}
\label{anexo:syncon}
% ------------------------------------------------------------------
Los \textit{Compensadores Síncronos} son máquinas rotativas síncronas
operadas en vacío ---sin turbina primaria--- que aportan inercia
rotacional genuina y capacidad de inyección de corrientes de falta de
300--400~\% de su valor nominal. Su papel en la transición hacia sistemas
dominados por IBR es el de proveer rigidez nodal (elevado $S_{sc}$) y
masa cinética natural, complementando la velocidad y precisión de los
inversores GFM. La estrategia \textit{Brownfield} aprovecha alternadores
de centrales clausuradas, reduciendo drásticamente costes de capital.

% ------------------------------------------------------------------
\section{Curva de Pato (\textit{Duck Curve})}
\label{anexo:duck_curve}
% ------------------------------------------------------------------
La \textit{curva de pato} describe el perfil diario de demanda neta de
regulación en sistemas con alta penetración solar: una depresión profunda
durante las horas centrales del día (cuando el consumo base es bajo pero
la generación solar es máxima) seguida de una rampa vespertina
pronunciada. La primavera es el período de máxima profundidad y
vulnerabilidad. En el caso del 28A, la profundidad del valle coincidió
con una rampa de inyección solar extraordinariamente aguda, dejando al
sistema con mínima capacidad de absorción de reactiva en el instante
crítico.

% ------------------------------------------------------------------
\section{Sistema VOLTAIRE}
\label{anexo:voltaire}
% ------------------------------------------------------------------
El sistema VOLTAIRE (integrado en el Centro de Control de Energías
Renovables, CECRE) es la arquitectura implantada por REE para la
regulación dinámica de tensión en el sistema peninsular. Opera en dos
capas jerárquicas: la Regulación Terciaria optimiza globalmente el
perfil de tensiones a nivel nacional mediante flujos de cargas óptimos
reactivos; la Regulación Secundaria envía consignas telemáticas en
tiempo real a los inversores de los parques renovables con una resolución
de muestreo máxima de 4~s, cerrando un bucle de control proporcional-integral
que mantiene la tensión dentro de márgenes fijos. La retribución del
servicio sigue la fórmula dinámica indexada al PMD establecida en la
revisión de junio de 2025 del P.O.~7.4.

% ------------------------------------------------------------------
\section{Centro de Control de Energías Renovables (CECRE)}
\label{anexo:cecre}
% ------------------------------------------------------------------
El \textit{Centro de Control de Energías Renovables} (CECRE) es la
entidad operativa de REE responsable de la monitorización y despacho en
tiempo real de los parques renovables y los sistemas de almacenamiento,
así como de la ejecución de los algoritmos de control de tensión a través
del sistema VOLTAIRE. Integra telemedida de alta resolución (PMU) de
instalaciones significativas y ejecuta subastas sub-horarias de servicios
ancilares.

% ------------------------------------------------------------------
\section{\textit{Power System Stabilizers} y \textit{Power Oscillation
Damping} (PSS/POD)}
\label{anexo:pss_pod}
% ------------------------------------------------------------------
Los \textit{Power System Stabilizers} (PSS) y los sistemas de
\textit{Power Oscillation Damping} (POD) son módulos de control
adicionales instalados en inversores (especialmente en modo GFM) que
inyectan señales contrafase diseñadas para amortiguar oscilaciones
electromecánicas de pequeña y gran perturbación. Detectan oscilaciones
de frecuencia o tensión y producen inyecciones o absorciones de potencia
sincronizadas que cancelan la amplitud oscilatoria, previniendo las
oscilaciones sub-síncronas (SSO) observadas durante el colapso del 28A.
El Real Decreto 997/2025 los hace obligatorios en inversores GFM de nueva
instalación con potencia superior a 10~MW.

% ------------------------------------------------------------------
\section{\textit{Headroom}: Reserva de Capacidad del Inversor}
\label{anexo:headroom}
% ------------------------------------------------------------------
El \textit{headroom} es la fracción de la capacidad aparente máxima
($S_{\max}$) que un inversor GFM debe mantener reservada sin utilizarla
para la inyección de potencia activa en estado estacionario. Esta reserva
es necesaria para garantizar que el inversor dispone de margen suficiente
para actuar ante perturbaciones rápidas de tensión o frecuencia. Exigir
\textit{headroom} reduce los ingresos del mercado de energía, lo que
constituye la fricción económica estructural que justifica la creación
de mercados de Servicios Esenciales de Confiabilidad (ERS) para
remunerar explícitamente esta capacidad de respuesta.

% ------------------------------------------------------------------
\section{Código de Red de Requisitos para Generadores (NC~RfG~2.0)}
\label{anexo:nc_rfg}
% ------------------------------------------------------------------
El \textit{Network Code on Requirements for Generators} (NC~RfG), ahora
en su versión 2.0, es el estándar técnico europeo que establece los
requisitos que deben cumplir las nuevas instalaciones de generación y
almacenamiento conectadas a la red. La versión 2.0, propuesta tras el
colapso ibérico, introduce la obligatoriedad de capacidades
\textit{grid-forming} para nuevas instalaciones de potencia significativa
(Tipos B, C, D), exigiendo que la tecnología habilite la inercia
sintética, el soporte dinámico de tensión y la inyección de corrientes
de cortocircuito robustas.

% ------------------------------------------------------------------
\section{\textit{Low Voltage Ride Through} (LVRT)}
\label{anexo:lvrt}
% ------------------------------------------------------------------
El \textit{Low Voltage Ride Through} (LVRT) es la capacidad de un
inversor para mantener la inyección de energía durante un hueco de
tensión en lugar de desconectarse por protección. Los requisitos del
LVRT en España están regulados por el P.O.~12.3 e incluyen el parámetro
dinámico $k$ (factor de proporcionalidad de corriente reactiva respecto
a la profundidad del hueco). El apagón del 28A evidenció que en redes
con $\mathrm{SCR} < 2$, la inyección masiva de reactiva según los
perfiles tradicionales de LVRT puede amplificar la inestabilidad en lugar
de contenerla, requiriendo revisión de la coordinación entre el control
LVRT y la debilidad de red.

